% !TeX root = ../tg.tex
\chapter{结论与展望}

本课题首次研究了标签噪声和非独立同分布数据同时存在时分割-联邦学习的精度与能耗情况。我们提供了关于标签噪声和非独立同分布数据如何影响分割-联邦学习收敛的第一个理论分析，考虑了均方误差损失和交叉熵损失。分析表明，在均方误差损失下，收敛速率的上界随着噪声方差的增大而增大，在交叉熵损失下，收敛速率的上界随着噪声率的增大而增大。为了解决这个问题，我们提出了一种基于分割-联邦学习的算法，可以有效地在非独立同分布数据下检测和校正噪声标签。在公共数据集和实际测试平台上的实验证明了该方法在准确性和能效方面的优越性能。

在未来的工作中，我们将重点研究分割-联邦学习与无反向传播训练方法的结合，在保证用户隐私的前提下最大程度地优化分割-联邦学习的精度与能耗。